\section{Problem Statement}
\label{sec:problem}

We consider programming a reusable \textit{nonprehensile manipulation} strategy from a single demonstration $D$ and a language instruction $l$. The demonstration is collected in one environment instance $E_0$, while the resulting program should solve a family of task variants $\mathcal{E}$ that share the same task objective but may differ in object geometry, pose, physical properties, and scene configuration.

We assume that no task-specific manipulation primitives, explicit success criterion, or interactive simulation environment are provided a priori. Given only $(D, l)$, the agent must first construct a set of executable manipulation primitives
\[
\mathcal{P} = {P_1,\ldots,P_N},
\]
and then compose these primitives into a strategy $\Pi$. The resulting program should capture the manipulation structure underlying the demonstration and generalize beyond the demonstrated instance:
\[
(D,l)
\longrightarrow
(\mathcal{P},\Pi),
\qquad
\Pi(E;\mathcal{P}) \text{ succeeds for } E \in \mathcal{E}.
\]
The goal is therefore not to reproduce the demonstrated motion, but to recover reusable primitives and a strategy that can adapt to unseen task variants.
